\section*{Сокращения}
\addcontentsline{toc}{section}{Сокращения}

В настоящем отчёте применены следующие сокращения:

\begin{description}
    \item[SLURM] Slurm Workload Manager, система диспетчеризации заданий и управления ресурсами вычислительного кластера. Исторически написание SLURM расшифровывалось как Simple Linux Utility for Resource Management.
    \item[MUNGE] MUNGE Uid 'N' Gid Emporium, служба аутентификации, используемая SLURM для проверки сообщений между узлами.
    \item[SSH] Secure Shell, протокол защищённого удалённого доступа.
    \item[NAT] Network Address Translation, механизм трансляции сетевых адресов.
    \item[UID] User Identifier, числовой идентификатор пользователя в Unix-подобных системах.
    \item[UTC] Coordinated Universal Time, всемирное координированное время.
    \item[IP] Internet Protocol, сетевой протокол, используемый для адресации узлов в сети.
\end{description}

\newpage
